\documentclass[letterpaper,titlepage,12pt,spanish]{article}
\usepackage[utf8]{inputenc}
\usepackage{times}
\usepackage{hyperref}
\usepackage[spanish]{babel}
\usepackage{graphics}
\usepackage{graphicx}
\usepackage[]{graphicx}
\usepackage{enumerate}
\usepackage{verbatim}
\usepackage{float}
\usepackage[spanish]{layout}
\usepackage{latexsym}
\usepackage[scriptsize,it]{caption}
\usepackage{fancyhdr}
\setlength{\headheight}{15pt}
\setlength{\headsep}{14pt} %OJO
\setlength{\footskip}{26pt} %OJO
\pagestyle{fancy}
\usepackage{colortbl}
\usepackage{multirow}
\hyphenation{EAFIT}
%\pagestyle{empty}
\sloppy
\hypersetup {bookmarksopen,bookmarksnumbered,colorlinks,linkcolor=blue,legalpaper,pdftitle=Libro de proyecto,pdfauthor=Carlos Urrego Moreno Juan David Pineda Cardenas Santiago Toro Acelas Héctor Espoz Alvarez}

\title{Centro de Computación Científica APOLO \\ Informe de Gestión 2013}
\vspace{1cm}
\author{Juan Guillermo Lalinde--Pulido\\jlalinde@eafit.edu.co \and Juan David Pineda--Cárdenas\\jpineda2@eafit.edu.co}
\date{
	\includegraphics[width=0.35\textwidth]{imgs/logo_apolo}\\[2cm]
	Universidad EAFIT\\
	Medellín, Colombia
}
\begin{document}
\fancyhead[]{}
\fancyhead[C]{Centro de Computación Científica Apolo}
\maketitle
%\begin{abstract}
%Este documento hace un breve resumen acerca del funcionamiento y mejoras realizados durante el año 2013 en el Centro de Computación Científica APOLO.
%\end{abstract}
\thispagestyle{empty}
\tableofcontents
\pagenumbering{arabic}
\newpage
\section{Objetivos}

	\begin{itemize}
		\item Construir un manual para la instalación de las librerías de Boost.
		\item Generar Documentación sobre el manejo de las librerías de Boost.
		\item Correr códigos de ejemplo para la ejecución de las librerías de Boost.
	\end{itemize}

\section{Alcance y productos}
	\begin{itemize}
		\item Descripción de las librerías de Boost.
		\item Introducción a las librerías de utilización de grafos instaladas en Apolo.
		\item Introducción al manejo de Grafos con Boost (representación conceptual).
	\end{itemize}

\section{Productos}
	\begin{itemize}
		\item Manual de Instalacion de librerias Para Boost.
		\item Ejemplo en código documentado y manual de como hacer un programa utilizando Boost para grafos.	
		\item 
	\end{itemize}

\section{Usabilidad}

	Utilizar los GPS de los taxis para hacer análisis de la malla vial de la ciudad. 
	Movilización y Rutas óptimas para el recorrido de Bomberos, Ambulancias, Policía para la atención de emergencias y desastres

	Movilidad, Seguridad, Educación, Tecnología.


\end{document}

